% title: Prime-free Recurrence Parameters
% short-title: Prime-free Recurrences
% abstract: We construct infinitely many recurrence parameters for which every term is non-prime, using a factorization of even Vieta--Fibonacci polynomial values.
% subjclass: 11B37, 11A41, 03B35
% keywords: recurrence sequence, non-prime term, Chebyshev polynomial, Lean 4
% author: Mario Hernández M.

\section{Proof of the theorem}

\begin{theorem}
There exist infinitely many positive integers \(x\) such that the sequence
\((a_n)_{n\ge 0}\) defined by
\[
a_0=1,\qquad a_1=x+1,\qquad a_{n+2}=xa_{n+1}-a_n
\]
has no prime term.
\end{theorem}

\begin{proof}
In fact, we prove the stronger statement that such parameters \(x\) are
arbitrarily large.

Let \(S_0(t)=1\), \(S_1(t)=t\), and
\[
S_{n+2}(t)=tS_{n+1}(t)-S_n(t).
\]
Choose any integer \(y\ge 3\) and put \(x=y^2-2\).  Define
\[
a_n=S_{2n}(y).
\]
The recurrence for the even-indexed subsequence gives
\[
S_{2n+4}(y)=(y^2-2)S_{2n+2}(y)-S_{2n}(y),
\]
so this is exactly a sequence with parameter \(x\).  Its first two values are
\[
a_0=1,\qquad a_1=S_2(y)=y^2-1=x+1.
\]

It remains to show that no \(a_n\) is prime.  The case \(n=0\) is \(a_0=1\).
For \(n\ge 1\), use the addition identity
\[
S_{m+n}(t)=S_m(t)S_n(t)-S_{m-1}(t)S_{n-1}(t)
\]
with \(m=n\).  This gives
\[
a_n=S_{2n}(y)=S_n(y)^2-S_{n-1}(y)^2
      =\bigl(S_n(y)+S_{n-1}(y)\bigr)
       \bigl(S_n(y)-S_{n-1}(y)\bigr).
\]
We prove by induction that \(S_n(y)>0\) and
\(S_{n+1}(y)-S_n(y)>1\).  The case \(n=0\) follows from
\(S_0(y)=1\), \(S_1(y)=y\), and \(y\ge3\).  If the assertion holds at \(n\),
then
\[
\begin{aligned}
S_{n+2}(y)-S_{n+1}(y)
&=(y-1)S_{n+1}(y)-S_n(y)\\
&=(y-2)S_{n+1}(y)+\bigl(S_{n+1}(y)-S_n(y)\bigr)>1,
\end{aligned}
\]
and \(S_{n+2}(y)>S_{n+1}(y)>0\).  Hence both displayed factors are strictly
greater than \(1\), so \(a_n\) is composite for \(n\ge1\), while \(a_0=1\) is
non-prime.

Taking \(y\) arbitrarily large gives arbitrarily large parameters
\(x=y^2-2\).  Thus there are infinitely many positive integers \(x\) with the
desired property.  The formal Lean statement
\lean{Biblioteca.Demonstrations.CompositeTermsInRecurrence.infinitely_many_positive_parameters_without_prime_terms}
records this as a theorem with arbitrarily large parameters; the constructed
sequence is represented over \(\mathbb Z\), and \texttt{natAbs} is used for the usual
natural-number primality test.
\end{proof}
